\documentclass[10pt,a4paper]{amsart}
\usepackage[margin=19mm]{geometry}
\usepackage{amsmath,amssymb,mathtools}
\usepackage[scr=rsfs,cal=euler]{mathalfa}
\usepackage{xfrac}
\usepackage[unicode,hidelinks,bookmarks=false]{hyperref}
\newcommand{\ie}{i.e.\ }
\newcommand{\cf}{cf.\ }
\newcommand{\resp}{respectively\ }
\newcommand{\Subsection}{\subsection}
\providecommand{\nobreakdash}{\nobreak\hbox{-}}
\setlength{\emergencystretch}{2em}
\multlinegap0pt
\usepackage[shortlabels]{enumitem}
\usepackage{amssymb}
\usepackage{mathtools}
\usepackage{tikz}
\usepackage{xcolor}
\newtheorem{thm}{Theorem}
\newtheorem{lemma}{Lemma}
\DeclareMathOperator{\AO}{AO}
\DeclareMathOperator{\ASp}{ASp}
\DeclareMathOperator{\AU}{AU}
\newcommand{\GL}{\mathrm{GL}} 
\DeclareMathOperator{\Or}{O}
\DeclareMathOperator{\Sp}{Sp}
\DeclareMathOperator{\U}{U}
\title{On the proportion of derangements \\in affine classical groups}
\author{Jessica Anzanello}
\date{Source: arXiv:2508.07093v2 (CC BY 4.0)}
\begin{document}
\maketitle

\noindent\emph{Source and licence.} On the proportion of derangements \\in affine classical groups. Version arXiv:2508.07093v2 (CC BY 4.0), distributed by arXiv under the Creative Commons Attribution 4.0 International licence, http://creativecommons.org/licenses/by/4.0/, as recorded in the arXiv licence metadata for this exact version and shown on the abstract page https://arxiv.org/abs/2508.07093v2. Full licence notice: \texttt{licenses/arxiv-2508.07093-CC-BY-4.0.txt}.\par\smallskip

\noindent\textit{Editorial scope.} Counted unit: the article's thm main (source0.tex lines 162--180). The article's complete original proof of it is delivered with it (source0.tex lines 563--591, one \texttt{proof} environment of 29 lines, which the article places away from the statement). No mathematics is written, repaired or re-derived: the statement and the proof are the article's own lines, in the article's own notation, and no retained line is substituted or re-flowed.  Retained uncounted as the source-local context the proof needs: lines 192--208; lines 342--356; lines 404--407; lines 534--560.  Nothing was omitted from any retained span, so the package carries no omission record.  No retained line contains a citation, so the wrapper carries no bibliography and no bibliography entry.  No retained line contains a figure environment, an image-inclusion command or drawing code, so no image is delivered and the package declares no asset dependency.  Editorial formatting only, in addition: an AMS \texttt{amsart} preamble that restates the article's own declarations and macros used by the retained lines, namely the environments \texttt{thm}, \texttt{lemma} on separate counters, together with the standard packages those lines need.  The retained environments print in this document as Theorem 1, Lemma 1 in order of appearance, whereas the article prints the same items with the numbering of its own full document.  The complete native source of the article as distributed in the arXiv e-print for this exact version is bundled unchanged as \texttt{sources/182988/source0.tex}.
\begin{thm}
\label{p_main}
The following formulas hold.
\begin{enumerate}
\item \label{p_Unitary}
$\displaystyle \delta_p(\AU_m(q))=
\frac{1}{q^m(q+1)}\sum_{i=1}^m\frac{(-1)^i((-q)^{i+1}-1)}{(-q)^{i(i+1)/2}(-1/q)_{m-i}};$
\item \label{p_Symplectic}
$\displaystyle \delta_p(\ASp_{2m}(q))=\frac{1}{q^m(q+1)}\sum_{i=1}^{m}\frac{(-1)^{i-1}(q^{2i+1}+1)}{q^{i(i+1)}(1/q^2)_{m-i}};$
\item \label{p_Orthogonal_2m+1}
$\displaystyle \delta_p(\AO_{2m+1}(q))=\frac{1}{2q^m(1/q^2)_{m}}+\frac{1}{2q^{m+1}}\sum_{i=0}^{m}\frac{(-1)^{i-1}}{q^{i(i+1)}(1/q^2)_{m-i}}$, with $p$ odd;
\item \label{p_Orthogonal_2modd}
$\displaystyle \delta_p(\AO^{\pm}_{2m}(q))=\frac{q^m\pm1}{2q^{2m}}\sum_{i=1}^{m}\frac{(-1)^{i-1}}{q^{i(i-1)}(1/q^2)_{m-i}}$, with $p$ odd;
\item \label{p_Orthogonal_2meven}
$\displaystyle \delta_2(\AO^{\pm}_{2m}(q))=\frac{1}{2q^{m-1}(1/q^2)_{m-1}}\pm \frac{1}{2q^{2m}}\sum_{i=1}^{m}\frac{(-1)^{i-1}}{q^{i(i-1)}(1/q^2)_{m-i}}.$
\end{enumerate}
\end{thm}

\begin{lemma}
\label{steinberg}
The following formulas hold.
\begin{enumerate}
    \item \label{p_u}
    $\displaystyle \Delta_u(\U_m(q))=\frac{1}{q^m(-1/q)_m}$;
    \item \label{p_sp} $\displaystyle \Delta_u(\Sp_{2m}(q))=\frac{1}{q^m(1/q^2)_m}$;
    \item \label{p_o_2m1} $\displaystyle \Delta_u(\Or_{2m+1}(q))=\frac{1}{2q^m(1/q^2)_m}$, if $p$ is odd;
    \item \label{p_o2m} $\displaystyle \Delta_u(\Or^{\pm}_{2m}(q))=
    \begin{cases}
    \displaystyle\frac{q^{m^2-m}}{2(q^m \mp 1)\prod_{i=1}^{m-1}(q^{2i}-1)}, & \text{ if } p \text{ is odd};\\
    \displaystyle\frac{q^{(m-1)^2}(q^m+q^{m-1}\mp1)}{2(q^m \mp 1)\prod_{i=1}^{m-1}(q^{2i}-1)}, & \text{ if } p=2.
    \end{cases}$
\end{enumerate}
\end{lemma}

\begin{equation}
\label{fact_Z_u}
Z_{\U}= \prod_{\phi \neq z, \phi=\widetilde{\phi}} \left( \sum_{\lambda}x_{\phi,\lambda}\frac{(-y)^{|\lambda|\deg(\phi)}}{c_{\GL,-q^{\deg(\phi)}}(\lambda)}\right)\cdot \prod_{\substack{\{\phi, \widetilde{\phi}\} \\ \phi \neq \widetilde{\phi}}} \left( \sum_{\lambda}x_{\phi,\lambda}x_{\widetilde{\phi},\lambda}\frac{y^{2|\lambda|\deg(\phi)}}{c_{\GL,q^{2\deg(\phi)}}(\lambda)}\right). 
\end{equation}

\begin{lemma}
\label{fact_D'_u}
We have
$$D'_{\U}=T^{-1}_{\U}(1-y)^{-1}U'_{\U}.$$
\begin{proof}
Recall that for a unipotent element $a \in \U_m(q)$, the partition $\lambda$ associated with $z-1$ satisfies $\lambda'_1=\dim \ker(a-1)$. It follows from the definition of the cycle index $Z_{\U}$ 
%that $u'_m(q)$ is the coefficient of $y^m$ in $Z_{\U}$ under the following substitutions:
%\begin{itemize}[label=$\cdot$]
  %  \item we set  $x_{\phi,\lambda}=0$ and $x_{\widetilde{\phi},\lambda}=0$ whenever $\phi\neq z-1$  and $\lambda \neq \emptyset$, 
   % \item we set $x_{\phi,\emptyset}=1$, 
   % \item and we set $x_{z-1,\lambda}=q^{-2\lambda'_1}$. 
%\end{itemize}
 %Similarly, the definition of cycle index yields 
 that $d'_m(q)$ is the coefficient of $y^m$ in $Z_{\U}$ when we assign the variables $x_{\phi,\lambda}$ and $x_{\widetilde{\phi},\lambda}$ the value 1 when $\phi \neq z-1$, and set $x_{z-1,\lambda}=q^{-2\lambda'_1}.$
 Moreover, observe that if all variables $x_{\phi,\lambda}$ and $x_{\widetilde{\phi},\lambda}$ are set to $1$ in $Z_{\U}$, then we obtain $(1-y)^{-1}$. Similarly, note that $u'_m(q)$ is the coefficient of $y^m$ in the factor 
  $$\sum_{\lambda}x_{z-1,\lambda}\frac{(-y)^{|\lambda|}}{(-q)^{\sum_{i}(\lambda'_i)^2}\prod_i(-1/q)_{m_i(\lambda)}}$$
of $Z_{\U}$, when we set the variables $x_{z-1,\lambda}$ equal to $q^{-2\lambda'_1}$. If we instead set these variables $x_{z-1,\lambda}$ equal to 1,
the coefficient of $y^m$ in $\sum_{\lambda}\frac{(-y)^{|\lambda|}}{(-q)^{\sum_{i}(\lambda'_i)^2}\prod_i(-1/q)_{m_i(\lambda)}}$ is equal to the proportion of unipotent elements in $\U_m(q)$, which, by Lemma \ref{steinberg}, is $\frac{1}{q^m(-1/q)_m}$.
Therefore, 
 $$\sum_{\lambda}\frac{(-y)^{|\lambda|}}{(-q)^{\sum_{i}(\lambda'_i)^2}\prod_i(-1/q)_{m_i(\lambda)}}=
 \sum_{m=0}^{\infty}\frac{y^m}{q^m(-1/q)_m}=T_{\U}.$$
Using the factorisation of $Z_{\U}$ in Eq. \eqref{fact_Z_u}, these observations together imply that
 \begin{equation*}
 (1-y)^{-1}=D'_{\U}\cdot U'^{-1}_{\U} \cdot T_{\U}. \qedhere 
\end{equation*}
\end{proof}
\end{lemma}

\begin{thm}
\label{main}
The following formulas hold.
\begin{enumerate}
\item
\label{Unitary} 
$\displaystyle \delta(\AU_m(q))=\frac{1}{q+1}\left(1-\frac{1}{(-q)^{m(m+3)/2}}\right);$
\item
%previously: conj_der
\label{Symplectic} 
$\displaystyle \delta(\ASp_{2m}(q))=\frac{1}{q+1}\left(1-\frac{1}{(-q)^{m(m+2)}}\right);$
\item
\label{Orthogonal_odd}
$\displaystyle \delta(\AO_{2m+1}(q))=\frac{1}{2}+ \frac{(-1)^{m-1}}{2q^{(m+1)^2}};$
\item
\label{Orthogonal_even}
$\displaystyle \delta(\AO^{\pm}_{2m}(q))=\frac{1}{2} \pm \frac{(-1)^{m-1}}{2q^{m(m+1)}}$.
\end{enumerate}
\end{thm}

\begin{proof}[Proof of Theorem  $\ref{main}\eqref{Unitary}$]
Applying Theorem \ref{p_main}\eqref{p_Unitary}, we first obtain:
\begin{align*}
u'_m(q)& =\frac{1}{q^m(-1/q)_m}-u_m(q) \\
& =\frac{1}{q^m(-1/q)_m}-\frac{1}{q^m(q+1)}\sum_{i=1}^m\frac{(-1)^{i}((-q)^{i+1}-1)}{(-q)^{i(i+1)/2}(-1/q)_{m-i}}\\
& = \frac{1}{q^m(q+1)}\sum_{i=0}^{m}\frac{(-1)^{i}(1-(-q)^{i+1})}{(-q)^{i(i+1)/2}(-1/q)_{m-i}}.
\end{align*}
We now observe that the following factorisation holds for $U'_{\U}$:
\begin{align}
    \nonumber
U'_{\U}=\sum_{m=0}^{\infty}u'_m(q)y^m&=\sum_{m=0}^{\infty}\frac{1}{q+1}\sum_{i=0}^{m}\frac{(-1)^i(1-(-q)^{i+1})}{(-q)^{m-i}(-q)^{i(i+3)/2}(-1/q)_{m-i}}(-y)^m\\
\nonumber
    &=\sum_{m=0}^{\infty}\frac{1}{q+1}\sum_{i=0}^{m}\frac{(-1)^i(1-(-q)^{i+1})}{(-q)^{i(i+3)/2}}(-y)^i\cdot\frac{1}{(-q)^{m-i}(-1/q)_{m-i}}(-y)^{m-i}\\
    \label{fattorizz}
    &=\overline{D}_{\U}\cdot T_{ \U},
\end{align}
where $\overline{D}_{\U}$ is defined as
$$\overline{D}_{\U}:=\frac{1}{q+1}\sum_{i=0}^{\infty}\frac{1-(-q)^{i+1}}{(-q)^{i(i+3)/2}}y^i.$$
By Eq. \eqref{fattorizz} and Lemma \ref{fact_D'_u}, we have 
\begin{align*}
D'_{\U}&=T^{-1}_{\U}\cdot U'_{\U} \cdot (1-y)^{-1}= \overline{D}_{\U} \cdot (1-y)^{-1} \\
&=\frac{1}{q+1}\sum_{i=0}^{\infty}\frac{1-(-q)^{i+1}}{(-q)^{i(i+3)/2}}y^i\cdot \sum_{j=0}^{\infty}y^j=\sum_{m=0}^{\infty}\frac{1}{q+1}\sum_{i=0}^{m}\frac{1-(-q)^{i+1}}{(-q)^{i(i+3)/2}}y^m.
\end{align*}
Hence,
$$d'_m(q)=\frac{1}{q+1}\sum_{i=0}^{m}\frac{1-(-q)^{i+1}}{(-q)^{i(i+3)/2}}$$
and
$$d_m(q)=1-d'_m(q)=\frac{1}{q+1}\sum_{i=1}^{m}\frac{(-q)^{i+1}-1}{(-q)^{i(i+3)/2}}=\frac{1}{q+1}\left(1-\frac{1}{(-q)^{m(m+3)/2}}\right),$$
where the last equality can be easily verified by induction on $m$.
\end{proof}

\end{document}
